\documentclass{article}
\usepackage{/globals/math}
\usepackage{/globals/de}
\begin{document}
\section{Ungleichungen}
Werden zwei Terme miteinander verglichen, aber nicht gleich gestellt, so wird eine Ungleichung aufgestellt. Hier wird ausgesagt, dass der eine Term kleiner, kleiner gleich, größer, oder größer gleich als der andere Term ist.

Es können die gleichen Äquivalenzumformungen genutzt werden, die für das Lösen von Gleichungen angewendet werden. Es ist nur wichtig zu beachten, dass bei einer Multiplikation oder einer Division mit einer negativen Zahl der Vergleichsoperator umgedreht wird.

\subsection{Lösungsmengen}
Die Lösungsmenge einer Ungleichung ist in der Regel nicht nur eine Wert oder ein paar Werte. Stattdessen gibt es oft unendlich viele Lösungen, alle Werte für eine Variable die eine bestimmte Kondition erfüllen, wie in einem Intervall zu sein. Die Schreibweise so einer Lösungsmenge ist 
\[
 \mL = \{x \in \R \dvert \text{Kondition}\}
\]
Die Kondition kann unter anderem ein einfacher Vergleich (\zb $x < x_1$), eine Eingrezung in ein Intervall (\zb $x_1 < x < x_2$) oder auch eine Kombination aus mehreren anderen Konditionen (\zb $x < x_1 \stext{oder} x > x_2$) sein.

\subsection{Nichtlineare Ungleichungen lösen}
Um nichtlineare Ungleichungen zu lösen müssen auf bekannte Art und Weise (Äquivalenzumformungen, p-q-Formel, \etc) die Schnittstellen beider Terme gefunden werden. Diese Stellen und die Intervalle zwischen ihnen können sich im Weiteren genauer angeguckt werden; durch einfaches Nachdenken oder Punktprobem kann bestimmt werden, ob sie Teil der Lösungsmenge sind oder nicht. 
\end{document}